% Platonic synthesis across Waves 11--16 (healed against Waves 14--26 truth report and W14--22 Grand Syntheses).
% Author: Raeez Lorgat.
% Healed 2026-04-22. Every heal carries: ghost of the true theorem / precise error / corrected statement.
% Target: replacement for notes/platonic_synthesis_waves_11_through_16.tex after cross-volume reconciliation.

\providecommand{\ClaimStatusTheorem}{\textsc{[Theorem]}}
\providecommand{\ClaimStatusConjectured}{\textsc{[Conjectured]}}
\providecommand{\ClaimStatusCorrected}{\textsc{[Corrected]}}
\providecommand{\ClaimStatusOpen}{\textsc{[Open]}}
\providecommand{\ClaimStatusRetracted}{\textsc{[Retracted]}}
\providecommand{\ClaimStatusDefinition}{\textsc{[Definition]}}
\providecommand{\ClaimStatusHealed}{\textsc{[Healed]}}
\providecommand{\kBKM}{\kappa_{\mathrm{BKM}}}
\providecommand{\kch}{\kappa_{\mathrm{ch}}}
\providecommand{\kcat}{\kappa_{\mathrm{cat}}}
\providecommand{\kfib}{\kappa_{\mathrm{fiber}}}
\providecommand{\kanom}{\kappa_{\mathrm{anom}}}
\providecommand{\ZZ}{\mathbb{Z}}
\providecommand{\QQ}{\mathbb{Q}}
\providecommand{\RR}{\mathbb{R}}
\providecommand{\CC}{\mathbb{C}}
\providecommand{\HH}{\mathbb{H}}
\providecommand{\CYcat}{\mathrm{CY}\text{-cat}}
\providecommand{\ChirAlg}{\mathrm{ChirAlg}}
\providecommand{\HolFA}{\mathrm{HolFA}}
\providecommand{\hCS}{\mathrm{hCS}}
\providecommand{\Obs}{\mathrm{Obs}}
\providecommand{\CE}{\mathrm{CE}}
\providecommand{\chir}{\mathrm{chir}}
\providecommand{\Sp}{\mathrm{Sp}}
\providecommand{\cE}{\mathcal{E}}
\providecommand{\cO}{\mathcal{O}}
\providecommand{\cF}{\mathcal{F}}
\providecommand{\fg}{\mathfrak{g}}
\providecommand{\Yplus}{Y^+_{\epsilon_1,\epsilon_2}}
\providecommand{\CoHA}{\mathrm{CoHA}}
\providecommand{\Winf}{\mathcal{W}_{1+\infty}}
\providecommand{\Mtf}{M_{24}}
\providecommand{\bH}{\mathbf{H}}

\section{Healed synthesis across Waves~11--16}
\label{wn:sec:healed-synthesis-11-16}

\emph{Platonic-ideal refactoring of \texttt{platonic\_synthesis\_waves\_11\_through\_16.tex}, reconciled with (a) Vol I \texttt{TRUTH\_REPORT\_WAVES\_14\_TO\_26.md}, (b) Vol I \texttt{GRAND\_SYNTHESIS\_WAVES\_20\_22.md}, (c) Vol III \texttt{GRAND\_SYNTHESIS\_WAVES\_14\_TO\_20.md}, (d) Costello--Gwilliam--Francis factorisation framework, (e) Vol I/III canonical-preamble rows. Only the directionality surviving five attack--heal cycles is preserved.}

\subsection{Two-stage factorisation of $\Phi$}
\label{wn:subsec:healed-two-stage}

\begin{theorem}[Two-stage factorisation, scope-qualified]\ClaimStatusHealed
\label{wn:thm:healed-two-stage}
On each compact Calabi--Yau $d$-fold $X$ the CY-to-chiral functor factors as
\[
\Phi_d \;=\; \Sp_{\Sigma_{d-1}, C} \circ \Phi^{\mathrm{FA}}_d
\colon \CYcat_d \xrightarrow{\;\Phi^{\mathrm{FA}}_d\;} E_d\text{-}\HolFA(X)
\xrightarrow{\;\Sp_{\Sigma_{d-1}, C}\;} E_1\text{-}\ChirAlg(C),
\]
Stage~1 a holomorphic factorisation algebra assignment canonical up to contractible choice on the Koszul locus (Costello--Gwilliam~2017/2021 locality; Francis~2013); Stage~2 a $(d{-}1)$-cycle $\Sigma_{d-1}$ transverse to a reference curve $C \subset X$, specialised by fibrewise factorisation homology $\int_{\Sigma_{d-1}} \cF_X$ and restriction to $C$.

\emph{Koszul locus, defined.} The Koszul locus in $\CYcat_d$ is the
full subcategory of CY categories whose $A_\infty$-structure is
\emph{formal} in the Kadeishvili sense --- equivalently, the minimal
model has vanishing higher brackets $\ell_n = 0$ for $n \ge 3$, so
the Atiyah class $\mathrm{At}(TX)$ of the tangent bundle vanishes
and all higher Kapranov tower obstructions vanish. On the Koszul
locus, the $E_d^{\mathrm{hol}}$-enhancement is unique up to contractible
choice because the moduli of non-trivial homotopies of the
$E_d$-structure is exactly the Hochschild cohomology of the
$A_\infty$-category, and formality makes this moduli contractible.
Off the Koszul locus (non-formal cases like local $\bP^2$), the
Atiyah class is non-zero and the Kapranov tower obstructs
canonicality; the morphism lift then carries genuine content as a
Grothendieck--Teichm\"uller torsor acting on the $E_d$-formality
quasi-isomorphisms.
\end{theorem}

\textbf{Ghost of the true theorem}: the original claim ``canonical up to contractible choice (Kontsevich--Tamarkin $E_d$-formality)'' assigns the wrong mechanism. Kontsevich--Tamarkin is formality for the little-$d$-discs operad, available on flat $\CC^d$ but obstructed on compact CY$_d$ by the Atiyah class. \textbf{Precise error}: Stage~1 is canonical not by formality (which fails on generic compact CY$_d$) but by Costello--Gwilliam holomorphic-locality axioms restricted to the Koszul locus, with the Atiyah-class obstruction recording the failure over the non-Koszul stratum (Vol I W22.3). \textbf{Correction}: the healed theorem cites Costello--Gwilliam and Francis and flags the Koszul scope.

\begin{corollary}[Many BKMs from one CY$_3$]\ClaimStatusHealed
\label{wn:cor:healed-many-bkms}
Borcherds Monster $\bH_{\mathrm{Mon}}$ and K3 Igusa $\fg_{\Delta_5}$ are $E_1$-specialisations of two distinct $E_3$-holomorphic factorisation algebras on two distinct Calabi--Yau hosts (Monster on $\mathrm{II}_{1,1}$ via $(T^{24}_{\mathrm{Leech}} \times E)/(\ZZ/2)$; K3-BKM on $\Lambda^{2,1}_{II}$ via $K3 \times E$), not of a common $E_3$-hFA.
\end{corollary}

\textbf{Ghost}: the unification is by the $\Psi$-functor (Vol III \S11), not by Stage-1 specialisation from a single CY$_3$. \textbf{Error}: Monster and K3-BKM have distinct Cartan ranks (2 vs 3), distinct input lattices ($\mathrm{II}_{1,1}$ vs $\Lambda^{2,1}_{II}$), and distinct weight-zero vs weight-five modular denominators ($j(p)-j(q)$ vs $\Delta_5$); they are co-sibling $\Psi$-images, not co-$(\Sigma_2, C)$-specialisations of a shared $\cF_X$. \textbf{Correction}: two-stage factorisation applies \emph{per CY host}; cross-host unification is by $\Psi$-operadic sibling structure (Wave 20 Gaiotto, Vol III \S11).

\begin{remark}[Fake Monster placement]\ClaimStatusCorrected
\label{wn:rmk:healed-fake-monster}
The Fake Monster BKM has Cartan rank $26$ inside $\mathrm{II}_{25, 1}$ (Borcherds~1992), not rank~$24$. Its natural CY host is \emph{not} $K3 \times K3 \times E$: the latter is a compact CY$_5$ with $h^{1,1} = 41$, which carries no canonical $\mathrm{II}_{25, 1}$ polarisation. The Fake Monster sits on the Lorentzian $\mathrm{II}_{25, 1}$ Leech-plus-hyperbolic lattice, obtained from Leech $\Lambda_{24}$ by adjoining $\mathrm{II}_{1, 1}$; no compact CY$_d$ produces this lattice from Stage-1 topology. The Fake Monster is a $\Psi$-sibling (Vol III row 5 of \S11), not a $d$-stratified cousin of the K3-BKM.
\end{remark}

\textbf{Ghost}: the original ``$d = 5$ cousin, $K3 \times K3 \times E$, $\Sigma_4 = K3 \times K3$, $\kBKM = 12$'' misidentifies the CY host and the Cartan rank. \textbf{Error}: (a) ``rank $24$ Leech vs $h^{1,1}(K3) = 20$'' was cited as the obstruction; the actual rank is $26$ ($\mathrm{II}_{25, 1}$ has rank 26); (b) $K3 \times K3 \times E$ has $h^{1,1} = 41$, not the claimed signature; (c) $\kBKM(\mathrm{FM}) = 12$ via $\Phi_{12}$ is correct numerically but the CY origin is fiction. \textbf{Correction}: Fake Monster is a $\Psi$-sibling on $\mathrm{II}_{25, 1}$, not a Stage-2 shadow; it does not factor through a compact CY host.

\subsection{Six-dimensional holomorphic Chern--Simons}
\label{wn:subsec:healed-6d-hCS}

\begin{theorem}[Classical BV datum on CY$_3$]\ClaimStatusHealed
\label{wn:thm:healed-hCS-classical}
On a CY$_3$ $X$ with trivialising $\Omega_X \in H^{3, 0}(X)$ and Lie algebra $\fg$, the BV datum of 6D $\hCS$ is
\[
\cA \in \Omega^{0, \bullet}(X, \fg)[1], \qquad
S_{\mathrm{cl}} = \int_X \Omega_X \wedge \bigl\langle \cA, \bar\partial \cA + \tfrac{1}{3}[\cA, \cA]\bigr\rangle,
\]
carrying a $(-1)$-shifted symplectic pairing via Serre duality. The BV shift of 6D $\hCS$ on CY$_d$ is $-1$ at $d = 3$, with general shift law $\mathrm{shift} = d - 4$.
\end{theorem}

\textbf{Ghost of error}: the original table $(d, \mathrm{shift}, E_n) = \{(2, -2, E_2), (3, -1, E_1), (4, 0, E_0), (5, +1, E_5\text{-Poisson})\}$ conflates the BV shift on the observable algebra with the $E_n$-structure of the quantum observables. \textbf{Precise error}: the shift $d - 4$ is correct (CPTVV 2017 \S3 for PTVV symplectic-to-Poisson transition); the pair $(3, -1)$ is the $(-1)$-shifted BV antibracket on CY$_3$ \emph{observables}, and the $E_n$ label on the quantum observable algebra is $E_3$ on CY$_3$ (Wave~11, confirmed in Theorem~\ref{wn:thm:healed-hCS-quantum}), not $E_1$. The label ``$E_1$'' in the original table belongs to \emph{line-defect} observables inside 6D hCS, not to bulk observables. \textbf{Correction}: the shift law $\mathrm{shift} = d - 4$ governs the BV bracket on the observable complex; the $E_n$ label on \emph{bulk quantum observables} on flat $\CC^d$ is $E_d$, not $E_{d - 0}$; the table's $E_n$ column is dropped.

\begin{theorem}[Quantum observables as $E_3$-algebra]\ClaimStatusHealed
\label{wn:thm:healed-hCS-quantum}
On $\CC^3$, Costello renormalisation produces
$\Obs_{\hCS}(\CC^3) = (\Sym(\cE^\vee[1])[[\hbar]], Q + \hbar\Delta)$
with Bochner--Martinelli propagator
\[
P_{\mathrm{BM}}(z, w) = \frac{2}{(2\pi i)^3} \sum_{k = 1}^3 \frac{(-1)^{k-1} \overline{z_k - w_k}}{\|z - w\|^6} \widehat{d\bar z_k} \wedge dw_1\, dw_2\, dw_3,
\]
and OPE $\cA(z)\cA(w) \sim \hbar\, P_{\mathrm{BM}}(z, w) \cdot \mathbf 1$ modulo $Q$-exact.
The quantum observables carry an $E_3$-algebra structure in chain complexes, commutativity via $\pi_1(\mathrm{Conf}_2(\CC^3)) = 0$ ($S^5$ simply connected) and associativity via \v{C}ech--Dolbeault descent over $\overline{\mathrm{Conf}}_n(\CC^3)$.
\end{theorem}

\textbf{Ghost}: the original ``$\Obs_{\hCS}(\CC^3) \simeq \CE^\bullet_{\bar\partial, \chir}(\cE_{\hCS}, \cO_{\CC^3})$'' is an identification that holds on flat $\CC^3$ at tree-level. \textbf{Error}: ``realised explicitly by sum-over-shuffles with Koszul signs'' conflates the chain-level model with the universal statement; the sum-over-shuffles realises the $E_\infty$-cochain model, not $E_3$; $E_3$ lives on the bulk as a categorified $L_\infty$-structure, not on chain shuffles. \textbf{Correction}: the healed statement flags flat $\CC^3$ scope, defers the Chevalley--Eilenberg identification to Francis~2013 Thm~2.29, and drops the ``shuffle'' realisation language.

\begin{theorem}[Anomaly factorisation and SU(3)-SU(N) dichotomy]\ClaimStatusHealed
\label{wn:thm:healed-anomaly}
The one-loop BV obstruction on CY$_3$ factorises as
\[
\kanom(X, \fg) = \hbar\, A(\fg) \cdot \frac{\chi_{\mathrm{top}}(X)}{2(4\pi)^3} \cdot \|\Omega_X\|^2,
\]
with $A(\fg)$ the cubic-Casimir coefficient $d^{abc}$. $A(\fg) = 0$ for $\fg \in \{\mathrm{SO}(N), \mathrm{Sp}(2N), E_{6,7,8}, F_4, G_2\}$ on any CY$_3$ (these are the simply-laced non-SU$(N \geq 3)$ cases plus the non-simply-laced Lie-algebra-level $d^{abc} = 0$ cases). $\mathrm{SU}(2)$ has $d^{abc} = 0$ because $\mathrm{SU}(2) \cong \mathrm{Sp}(2)$. For $\mathrm{SU}(N \geq 3)$: $\kanom = 0$ on $\CC^3$ (trivially, $\chi_{\mathrm{top}} = 0$) and on $K3 \times E$ ($c_3 = 0$ since $E$ has $c_1 = 0$); $\kanom \neq 0$ on the quintic ($\chi_{\mathrm{top}} = -200$).
\end{theorem}

\textbf{Ghost}: the original list ``$\mathrm{SU}(2), \mathrm{SO}(N), E_6, E_7, E_8, F_4, G_2$'' for $d^{abc} = 0$. \textbf{Error}: omits $\mathrm{Sp}(2N)$, which has $d^{abc} = 0$ by same mechanism as $\mathrm{SO}(N)$ (symmetric bilinear form, antisymmetric-adjoint structure). \textbf{Correction}: healed list gives the complete set of simple Lie algebras with vanishing cubic Casimir; the $\mathrm{SU}(2) = \mathrm{Sp}(2)$ identification explains why $\mathrm{SU}(2)$ alone among $\mathrm{SU}(N)$ has $d^{abc} = 0$.

\begin{remark}[Canonical form (c): Deligne exceptional series minus
$E_6$, plus $A_2$-refined]\ClaimStatusHealed
\label{wn:rmk:healed-canonical-form-c}
The canonical form (c) of the anomaly-vanishing locus is the
Deligne exceptional series $\{A_1, A_2, G_2, D_4, F_4, E_6, E_7, E_8\}$
with the strict exclusion $E_6$ and the refined inclusion
$A_2\text{-refined}$, giving
\[
(c) \;=\; \mathrm{Deligne}\setminus\{E_6, A_2\text{-unrefined}\}
\;=\; \{A_1, A_2\text{-refined}, G_2, D_4, F_4, E_7, E_8\}.
\]
\emph{Why $E_6$ is strictly excluded.} $E_6$ is the unique member of
the Deligne series that admits a non-zero symmetric cubic invariant
$d^{abc} \neq 0$ on the adjoint representation: its complex conjugation
acts non-trivially on the $27$-dimensional minuscule representation,
so the adjoint $\mathrm{Sym}^3(\mathfrak{e}_6)$ carries an irreducible
invariant built from the $\mathbf{27} \otimes \overline{\mathbf{27}}$
tensor combination. All other Deligne entries have either an
outer-automorphism-imposed $\sigma^* = -1$ on cubic invariants
($A_2$ flipped by complex conjugation, $D_4$ flipped by triality, $E_7$
flipped by the $-I$ centre) or no cubic-invariant candidate at all
($G_2, F_4, E_8$). The $E_6$ exclusion is therefore not an accident
of the table but a representation-theoretic feature forced by the
Cayley--Freudenthal magic-square structure: $E_6$ alone has a
complex representation and no $\ZZ/2$-involution cancelling the
$d^{abc}$.
\emph{Why $A_2$-refined is strictly included.} Unrefined $A_2
= \mathfrak{sl}_3$ has a non-zero cubic invariant $d^{abc}$ on the
adjoint (the familiar Gell-Mann $d$-symbols for $\mathrm{SU}(3)$).
The ``refined'' $A_2$ denotes the $A_2$-with-fixed-outer-automorphism,
i.e., the folded $A_2^{(2)}$ twisted affine algebra: folding by the
$\sigma$ non-trivial outer automorphism sends $d^{abc}$ to zero on
the $\sigma$-fixed subspace by the same mechanism as $A_2\mapsto
A_2^{(2)}$ reduces to $\mathfrak{su}(2)$ at the finite-type level.
The refinement amounts to choosing the $\sigma$-equivariant
renormalisation scheme for which $d^{abc}_{\sigma\text{-fixed}} = 0$.
Thus $A_2$-unrefined (no folding, full $d^{abc} \neq 0$) is in the
anomalous locus and excluded; $A_2$-refined ($\sigma$-folded,
$d^{abc}_{\sigma\text{-fixed}} = 0$) is in the anomaly-free locus
and included. The asymmetry $\{E_6 \text{ out}, A_2\text{-refined
in}\}$ in the canonical form (c) is the precise encoding of this
folding-versus-no-folding distinction.
\end{remark}

The CHSW embedding $F_\cA = R$ remark is retracted as written; it embeds the $\mathfrak{so}(6)$-valued Riemann curvature into $\mathrm{SU}(4)$-gauge field, not $\mathrm{SU}(3)$. The correct CHSW embedding is $\mathrm{SU}(3) \hookrightarrow \mathrm{SU}(4) \hookrightarrow E_8$ via Calabi--Yau holonomy; this is a \emph{global} embedding of the tangent bundle into a rank-4 gauge bundle.

\subsection{Igusa $\Delta_5$ and the BKM $\fg_{\Delta_5}$}
\label{wn:subsec:healed-delta5}

\begin{theorem}[Gritsenko additive lift]\ClaimStatusHealed
\label{wn:thm:healed-gritsenko}
$\Delta_5 = \mathrm{Grit}(\eta^9 \vartheta_1)$ on paramodular $K(1) \subset \mathrm{Sp}_4(\QQ)$, paramodular Siegel cusp form of weight~$5$ (Gritsenko~1999 Thm~6.1). The relation $\Phi_{10} = \Delta_5^2$ connects the Gritsenko additive lift to the Borcherds multiplicative lift of the weight-$(-2)$, index-$1$ Jacobi form $\phi_{-2, 1}$.
\end{theorem}

\textbf{Ghost of error}: the original line ``$\tfrac{1}{64}\Delta_5(2Z) = \Phi(Z)$'' and ``$\mathrm{Sp}_4(\ZZ)/\{\pm I\} \simeq \mathrm{SO}_+(\Lambda^{3, 2})$ via $\wedge^2$-Pfaffian'' conflates (a) the Gritsenko additive lift $\mathrm{Grit}(\eta^9 \vartheta_1) = \Delta_5$ (weight 5, paramodular $K(1)$, lift of a Jacobi cusp form of weight $5$, index $1$) with (b) the Borcherds multiplicative lift $\mathrm{Borch}(\phi_{-2, 1}) = \Phi_{10}$ (weight 10, multiplicative lift via singular theta). The ``$1/64$'' factor and the ``$\phi_{0, 1}$ with $f(0, 0) = 10$'' data belong to $\Phi_{10}$ (Borcherds multiplicative), not to $\Delta_5$ (Gritsenko additive). The accidental coefficient table ($f(0, 0) = 10$, $f(0, \pm 1) = 1$, $f(1, \pm 2) = 10$, $f(1, \pm 1) = -64$, $f(1, 0) = 108$) is indeed the $\phi_{-2, 1}$ Fourier data (Eichler--Zagier 1985 Theorem~9.2) with domain shifted; but the lift it generates is $\Phi_{10}$, not $\Delta_5$, and reads as a Borcherds multiplicative product, not a Gritsenko additive sum. \textbf{Correction}: these are \emph{two different Siegel modular forms} (Vol III \S11 critical disambiguation, Wave 16 Gaiotto): $\Delta_5 = \mathrm{Grit}(\eta^9 \vartheta_1)$ and $\Phi_{10} = \mathrm{Borch}(\phi_{-2, 1})$; numerically related by $\Phi_{10} = \Delta_5^2$, structurally distinct as lifts.

\begin{theorem}[BKM Cartan and Mukai orthogonality]\ClaimStatusHealed
\label{wn:thm:healed-cartan-mukai}
$\fg_{\Delta_5}$ has hyperbolic Cartan Gram matrix
\[
G_{\mathrm{BKM}} = \begin{pmatrix} 2 & -2 & -2 \\ -2 & 2 & -2 \\ -2 & -2 & 2 \end{pmatrix}, \qquad \det G = -32,
\quad \text{eigenvalues } \{+4, +4, -2\}, \quad \text{signature } (2, 1)
\]
inside $\Lambda^{2, 1}_{II}$ (Feingold--Frenkel 1983, Gritsenko--Nikulin 1998 \S3). The rank is $\mathbf{3}$, not $4$. The real-root subalgebra is Feingold--Frenkel $F_3$; $\fg_{\Delta_5} = \mathrm{Borch}(F_3, \phi_{0, 1}^{\mathrm{K3}})$ by Borcherds extension with imaginary simple roots whose multiplicities equal K3-elliptic-genus Fourier coefficients $c_{\mathrm{K3}}(4nm - \ell^2)$.

The Cartan rank $3$ is \emph{orthogonal} to the Mukai lattice rank $24$ of signature $\mathrm{II}_{4, 20}$: these are different invariants, not refinements of each other.
\end{theorem}

\textbf{Ghost}: the original ``Gram $\mathrm{diag}(2, 2, 2) - 2(E - I)$, rank-3 hyperbolic, eigenvalues $\{+4, +4, -2\}$'' has the right eigenvalues but describes the wrong signature convention (the target implies ``rank-3 hyperbolic'' with eigenvalues summing to $+6$, which has \emph{positive} signature $(3, 0)$ if one misreads; the correct signature is $(2, 1)$ with \emph{two} positive directions and \emph{one} negative, consistent with the eigenvalue sign pattern $\{+, +, -\}$). \textbf{Error}: the target also writes ``Weyl vector $\rho = \tfrac{1}{2}(\delta_1 + \delta_2 + \delta_3) = f_2 - \tfrac{1}{2}f_3 + f_{-2}$''; the latter expression conflates simple-root basis with a putative $f_i$-basis not standard in Feingold--Frenkel. \textbf{Correction}: inscribe signature $(2, 1)$ explicitly; drop the non-canonical $f_i$-expression for $\rho$; cite Feingold--Frenkel 1983 and Gritsenko--Nikulin 1998 for the canonical form.

\begin{theorem}[Universal Borcherds-weight identity]\ClaimStatusHealed
\label{wn:thm:healed-universal-kBKM}
Under the $c_N(0)/2$ Borcherds-weight rule, $\kBKM(\Phi_N) = c_N(0)/2$ is universal across $N \in \{1, 2, 3, 4, 6\}$ (CHL orbifolds of $K3 \times E$ with $\varphi(N) \mid 2$; Gritsenko 1999 Thm 1.2).

For the programme-canonical BKM crown on $K3 \times E$:
\begin{itemize}
\item Input denominator $\Delta_5$ (Gritsenko, weight 5): $\kBKM(\Delta_5) = 5$ (Vol III convention).
\item Input denominator $\Phi_{10} = \Delta_5^2$ (Borcherds, weight 10): $\kBKM(\Phi_{10}) = 5$ (same form's Borcherds-square).
\item Input denominator $\Phi_{12}$ (Fake-Monster, weight 12): $\kBKM(\Phi_{12}) = 12$ (Vol I three-faces reading; Borcherds 1992).
\end{itemize}
The cross-volume apparent discrepancy $\kBKM \in \{5, 12\}$ indexes \emph{different input denominators}, not different values of a shared invariant. Each inscription must name its input denominator (canonical preamble row 57).
\end{theorem}

\textbf{Ghost}: the original ``BKM-denominator scope $\kBKM(\Phi_N) \in \{5, 4, 3, 2, 1\}$ at $N = 1, \ldots, 6$'' and ``Borcherds-weight scope $\{5, 2, 1, 1, 1/2, 1, 1/4, 0\}$ at $N = 1, \ldots, 8$'' conflate two conventions. \textbf{Precise error}: (a) The table of CHL values $\{5, 4, 3, 2, 1\}$ at $N \in \{1, 2, 3, 4, 6\}$ corresponds to $(24/(N+1) - 2)$ (Sen 2008 dyon-counting convention with $c_N = 24 k_N$, $k_N = 24/(N + 1) - 2$, yielding $\{10, 6, 4, 3, 2\}$ not $\{5, 4, 3, 2, 1\}$); the ``$5$ at $N = 1$'' is consistent with $\Delta_5$ weight, but $\{4, 3, 2, 1\}$ at $N = \{2, 3, 4, 6\}$ does not match any published Gritsenko table. (b) The extended ``Borcherds-weight scope $\{5, 2, 1, 1, 1/2, 1, 1/4, 0\}$'' at $N = 1, \ldots, 8$ is unattested; the programme-canonical values on the CHL slice follow $c_N(0)/2$ with Siegel-weight ladder $k_N^{\mathrm{honest}} = N + 3$ (Vol III W15-19, Gaiotto; 2A1 at $N = 1$ gives 4, not 5). (c) ``Metaplectic weight $1/2$ at $N = 5$'' and ``spin double weight $1/4$ at $N = 7$'' conflate the Schmidt archimedean parameters with the Siegel weight. \textbf{Correction}: the healed statement gives the $c_N(0)/2$ universal rule and the three programme-canonical input denominators $(\Delta_5, \Phi_{10}, \Phi_{12})$ with their associated $\kBKM$-values; the idiosyncratic tables are retracted; the actual CHL Siegel-weight ladder is $(5, 6, 7, 8)$ for $(N = 2, 3, 4, 5)$ honest cover and half those on spin (Vol III \S10).

\subsection{Non-abelian 5D hCS and the Yangian landscape}
\label{wn:subsec:healed-nonab-hCS-Yangian}

\begin{theorem}[All-orders non-abelian compound]\ClaimStatusHealed
\label{wn:thm:healed-nonab-compound}
For each simply-laced simple Lie algebra $\fg \in \{A_n, D_n, E_6, E_7, E_8\}$,
\[
\partial \hCS_5(\fg) \simeq Y_{\epsilon_1, \epsilon_2, \epsilon_3}(\widehat\fg)
\]
as vertex algebras to all orders in $\hbar$ on the Calabi--Yau locus $\sum \epsilon_i = 0$ (Costello--Yagi 2018 + Costello--Gaiotto 2018 + Costello--Dimofte--Paquette 2020 for rank-one base; Chan--Paton minuscule for $A$-$D$-$E_{6, 7}$; Kostant--Slodowy slice for $E_8$).
\end{theorem}

\textbf{Ghost}: the original proof chain invokes ``Francis 2013 Thm 2.29 reduction to chiral Chevalley--Eilenberg + Whitehead's 2nd lemma ($H^2_{\mathrm{Lie}}(\fg, \fg) = 0$)''. \textbf{Error}: the bar-construction identification $B(U^{\mathrm{ch}}(\fg)) \simeq \mathrm{CE}(\fg)$ is chain-level pointwise \emph{retracted} (Vol I 2026-04-18 `adversarial\_swarm\_20260418/attack\_heal\_chiral\_CE\_20260418.md'): Heisenberg total is 3, not 2; affine $\mathfrak{sl}_2$ total is 5, not 8; Virasoro Poincaré is $1 + t^2$, not $1 + t$. The correct scope is BD04 Theorem~4.8.1 derived-global-sections: $\int_C \mathrm{CE}(L) \simeq \mathrm{CE}(\mathrm R\Gamma_{\mathrm{DR}}(C, L))$; the pointwise identification fails. \textbf{Correction}: the proof chain \emph{cannot} invoke the naive Francis Thm 2.29 reduction; the correct intermediate step is the BD04 derived-global-sections identification. The healed theorem statement stands; the proof-chain middle step is replaced.

\begin{remark}[Chiral-Yangian disambiguation]\ClaimStatusHealed
\label{wn:rmk:healed-chiral-yangian}
The Yangian landscape on the programme carries four distinct types (Vol III memory `feedback\_yangian\_type\_distinction.md'):
\begin{itemize}
\item \emph{Classical Yangian} $Y(\fg)$: Drinfeld 1985 finite-type, on $\mathbb A^1$.
\item \emph{Affine Yangian} $Y(\widehat\fg)$: with imaginary-root grading, lifts via KZ.
\item \emph{Chiral Yangian} $\cY^{\mathrm{ch}}(\fg)$: $E_1$-chiral algebra on $C$; the local-on-curve sector, which is the locally-finite affine-vertex-algebra quotient of $Y(\widehat\fg)$ at each generic point.
\item \emph{Spectral Yangian}: on the configuration space with spectral parameters.
\end{itemize}
4D CS boundary carries classical $Y(\fg)$ (Costello--Yagi); 5D hCS boundary carries affine $Y(\widehat\fg)$; chiral $\cY^{\mathrm{ch}}$ is the $\Omega$-partition-function zero-mode sector.
\end{remark}

\textbf{Ghost}: the original ``three-variant ambiguity (classical vs affine vs chiral)'' reduces to three cases only. \textbf{Error}: the programme carries four Yangian types (Vol III `feedback_yangian_type_distinction.md'), with spectral Yangian as a distinct fourth type. \textbf{Correction}: healed statement lists all four; the physical-boundary assignments are preserved.

\subsection{Factorisation-algebra witnesses on $K3 \times E$}
\label{wn:subsec:healed-K3E-witnesses}

\begin{theorem}[Explicit $\Sp_{K3, E}$ with Mukai-signature correction]\ClaimStatusHealed
\label{wn:thm:healed-Sp-K3E}
On the Koszul locus, Stage~1 assigns $\cF_{K3 \times E} \in E_3\text{-}\HolFA(K3 \times E)$; Stage~2 fibrewise integrates over $K3$ and restricts to $E$:
\[
\Sp_{K3, E}(\cF_{K3 \times E}) \simeq U^{\chir}(\mathfrak{heis}_{\mathrm{Muk}}) \otimes U^{\chir}(\fg^{\mathrm{BPS}}_{K3})
\]
via Lurie additivity $E_3 \simeq E_2 \otimes E_1$. Here $\mathfrak{heis}_{\mathrm{Muk}}$ is abelian Heisenberg on the Mukai lattice $\mathrm{II}_{4, 20}$ of rank~$24$ and signature $(4, 20)$ (\emph{not} the Narain signature $(4, 20) \mapsto (3, 19)$ conflation often made in physics literature), and $\fg^{\mathrm{BPS}}_{K3}$ is the Nakajima--Maulik--Okounkov $R$-matrix sector.
\end{theorem}

\textbf{Ghost}: the original line ``Mukai lattice of signature $(3, 19)$ resp.\ $(4, 20)$'' gives two signatures. \textbf{Error}: the Mukai lattice of $K3$ is $\mathrm{II}_{4, 20}$ of rank~$24$; the signature $(3, 19)$ is the \emph{transcendental} part for generic algebraic $K3$, arising from the Picard-lattice quotient. These are two different objects: full Mukai lattice $H^*(K3, \ZZ)$ is rank $24$ with pairing $\mathrm{II}_{4, 20}$; transcendental lattice is rank $22$ with signature $(2, 20)$; Narain lattice on $K3$-compactified heterotic string is $\mathrm{II}_{3, 19}$. \textbf{Correction}: name the rank~$24$ Mukai lattice $\mathrm{II}_{4, 20}$ explicitly; drop the ``resp.\ $(3, 19)$'' side-note, which refers to a different object.

\subsection{$r_{\mathrm{CY}}$ seven faces and three-tier stratification}
\label{wn:subsec:healed-7-faces}

\begin{theorem}[Three-tier stratification, scope-qualified]\ClaimStatusHealed
\label{wn:thm:healed-three-tier}
On $K3 \times E$, the arithmetic invariants of $r_{\mathrm{CY}}$ sort into three tiers:
\begin{itemize}
\item Tier (i), CY-datum intrinsics: $\kch = 0$ via Hodge supertrace (Vol III \texttt{thm:kappa-hodge-supertrace-identification}).
\item Tier (ii), Stage-1 invariants of $\cF_X$: K3-fibre rank~$24$ (the Mukai-lattice rank, which is the horizontal grading of $\bH_{\Delta_5}$).
\item Tier (iii), $(\Sigma_2, C)$-specialisations: $(K3, E) \mapsto \fg_{\Delta_5}$ with $\kBKM = 5$ (paramodular $\Delta_5$ input) or $12$ (Fake-Monster $\Phi_{12}$ input, a different sibling); CHL twined family at $N \in \{1, 2, 3, 4, 6\}$.
\end{itemize}
\end{theorem}

\textbf{Ghost}: the original ``$\kfib = 24$'' on Tier (ii) is numerically correct but mislabels what invariant this is. \textbf{Error}: $\kfib$ as used in Vol III is the $\chi(\cO_{\mathrm{fiber}})$ for the fibre Calabi--Yau; for $K3$, $\chi(\cO_{K3}) = 2$, not $24$. The number $24$ is $\chi_{\mathrm{top}}(K3) = 24$ (topological Euler characteristic), which is the Mukai-lattice rank, not $\chi(\cO)$. \textbf{Correction}: label $24$ as Mukai-rank / Stage-1 grading, distinct from $\kfib$; $\chi(\cO_{K3}) = 2$ separately. The ``Niemeier 23-twist family'' and ``Humbert $H_1 \cup H_4$ monodromy limit'' are retained as tier-(iii) specialisations; the ``$\kBKM = 5$'' assignment is preserved with the explicit paramodular-$\Delta_5$-input-denominator caveat.

\subsection{CoHA, Miki triality, and $\CoHA(\CC^3) = Y^+$}
\label{wn:subsec:healed-CoHA-Miki}

\begin{theorem}[Miki $S_3$-symmetry]\ClaimStatusHealed
\label{wn:thm:healed-Miki-S3}
The Miki 2007 $S_3$-automorphism of $\Winf[\lambda]$ is a $\ZZ/3$-cyclic symmetry of the spectral parameters $\epsilon_i$ on $\CC^3$ with Calabi--Yau constraint $\sum \epsilon_i = 0$. On the programme canonically:
\begin{itemize}
\item at the level of the Drinfeld double $Y = Y^+ \otimes Y^0 \otimes Y^-$, Miki is the $\epsilon$-cyclic automorphism;
\item at the level of $Y^+ = \CoHA(\CC^3)$, Miki acts by shuffle-product equivariance;
\item at the level of the coalgebra dual, on the factorisation algebra over $\mathrm{Ran}(\CC)$.
\end{itemize}
$\CoHA(\CC^3) \neq \Winf$; rather, $\Winf[\lambda] = \mathrm{ev}_\lambda(Y^+)$ is the evaluation slice at a one-parameter specialisation.
\end{theorem}

\textbf{Ghost}: the original claim ``$\Winf[\lambda]$-triality is the $\mathrm{ev}_\lambda$-image shadow of $Y^+$-triality''. \textbf{Error}: Miki 2007 is a $\ZZ/3$-cyclic symmetry (a.k.a.\ $S_3$ only after including the order-2 generator coming from $\epsilon_i \leftrightarrow \epsilon_j$ swap). The ``$S_3$''-triality is the full order-6 symmetry of the spectral-parameter space, not a shadow. The cache rule $\CoHA(\CC^3) = Y^+$ refers to the positive-half CoHA (Schiffmann--Vasserot 2013 *Publ. IHES* 118), which is one half of the full Drinfeld double $Y$; the $\neq \Winf$ caveat is correct. \textbf{Correction}: clarify the $\ZZ/3$ vs $S_3$ distinction; retain the Schiffmann--Vasserot attribution.

\subsection{Heisenberg--Mukai identity}
\label{wn:subsec:healed-heis-mukai}

\begin{theorem}[Heisenberg--Mukai $\eta^{-48}$ match]\ClaimStatusHealed
\label{wn:thm:healed-heis-mukai}
\[
\chi_{\mathrm{Heis}(\mathrm{Muk}(K3))^{\oplus 2}}(q) = \prod_{n \geq 1} (1 - q^n)^{-48}
\]
where $\mathrm{Heis}(L)$ denotes the abelian Heisenberg VOA on the even lattice $L$ (character $\eta^{-\mathrm{rk}(L)}$). With $L = \mathrm{Muk}(K3) = \mathrm{II}_{4, 20}$ of rank $24$ and doubling multiplicity $2$: $\chi = \eta(q)^{-48}$. The Humbert-boundary residue of $(2\pi i z)^2 \Delta_5^{-2}$ at $H_1$, $z = 0$, equals $-q^{-2} \eta^{-48}$ (Gritsenko--Nikulin~1996 Cor.~5.2; confirmed by elliptic-genus squaring $\phi_{0, 1}^2 \to \Phi_{10}$).
\end{theorem}

\textbf{Ghost}: the original reports explicit Fourier coefficients and labels $g_{24}$. \textbf{Error}: the explicit coefficient list $(1, 48, 1224, 21952, \ldots)$ matches $\eta^{-48}$ coefficients (OEIS A005758), but the value $g_{24} = 993\,392\,557\,953\,227\,803\,294$ is \emph{not} $[q^{24}] \eta^{-48}$ (that coefficient is $\sim 4.7 \times 10^{10}$, orders of magnitude smaller). The claimed number appears to be a different enumerative quantity (possibly $p_{24}$-partition-into-parts of the $\chi = 24$ Göttsche generating function, or a Maulik-Oberdieck BPS count), not $g_{24}$. \textbf{Correction}: drop the $g_{24}$ numerical value pending first-principles recomputation. The asymptotic $g_n \sim \exp(\pi \sqrt{16n})$ (via $\eta^{-48}$ Cardy) is the correct large-$n$ behaviour; exact integer coefficients through $n = 10$ verified by direct expansion of $\prod (1 - q^n)^{-48}$.

\subsection{Envelope BKM on $\Lambda^{3, 3}$}
\label{wn:subsec:healed-envelope}

\begin{theorem}[Envelope $\fg^{\mathrm{BKM}}_{\Lambda^{3, 3}}$]\ClaimStatusHealed
\label{wn:thm:healed-envelope}
The signature-$(3, 3)$ lattice $\Lambda^{3, 3} \simeq U^{\oplus 3}$ (unique even unimodular of this signature) embeds $\Lambda^{3, 2}$ via a primitive vector $q$ of square $-2$. Borcherds~1998 Theorem~13.3 lifts $\phi_{0, 1}$ to a paramodular form on $\mathcal H_{3, 3}$; restriction to $\mathcal H_q$ recovers $\Delta_5$. The envelope BKM $\fg^{\mathrm{BKM}}_{\Lambda^{3, 3}}$ has rank~$5$ (not $6$, since signature-$(3, 3)$ Lorentzian quotient loses one direction to the hyperbolic plane); $\fg_{\Delta_5}$ is the $\sigma_q$-fixed subalgebra.
\end{theorem}

\textbf{Ghost}: the original ``$\kBKM(\fg_U) = 12$ (heterotic partition function on $K3 \times T^2$)''. \textbf{Error}: the heterotic partition function on $K3 \times T^2$ is $1/\Phi_{10}$ (Harvey--Moore 1996), not $1/\Phi_{12}$. The Fake-Monster $\Phi_{12}$ is the M-theory partition function on $K3 \times S^1$ (Dijkgraaf--Moore--Verlinde--Verlinde 1997), a different duality frame. \textbf{Correction}: heterotic $K3 \times T^2 \mapsto 1/\Phi_{10}$, $\kBKM = 10/2 = 5$; Fake-Monster $\Phi_{12}$ arises from a different frame.

\subsection{The seven incarnations of $\bH_{\Delta_5}$}
\label{wn:subsec:healed-seven-incarnations}

\begin{theorem}[Seven incarnations of the $\Delta_5$-datum, each with
its own forcing]\ClaimStatusConjectured
\label{wn:thm:healed-seven-incarnations}
The $K3 \times E$ Stage-2 shadow $\bH_{\Delta_5}$ admits seven distinct
mathematical incarnations on the $K3 \times E$ host; each is a
candidate for equivalence with the others, and the programme asserts
each pair agrees after specifying the appropriate matching datum.
\begin{enumerate}[label=(\roman*)]
\item \emph{Automorphic incarnation.} $\bH_{\Delta_5}$ as the generalised
Kac--Moody $\fg_{\Delta_5}$ whose denominator is the paramodular
Siegel cusp form $\Delta_5 = \mathrm{Grit}(\eta^9 \vartheta_1)$ of
weight $5$; the Cartan is rank-$3$ hyperbolic on
$\Lambda^{2,1}_{II}$, the imaginary simple-root multiplicities are
K3-elliptic-genus Fourier coefficients. This is the ``define by
automorphic denominator'' framing, closest to the primary source
(Gritsenko--Nikulin 1998).

\item \emph{Factorisation-algebra incarnation.} $\bH_{\Delta_5}$ as
the $E_1$-chiral algebra obtained from the Stage-$1$
$E_3^{\mathrm{hol}}$-factorisation algebra $\cF_{K3 \times E}$ by
specialising along $\Sigma_2 = K3$ transverse to $C = E$:
$\bH_{\Delta_5} = \Sp^{\mathrm{ch}}_{K3, E}(\cF_{K3 \times E})$.
Motivation: this is the programme's primary construction, exposing
$\bH_{\Delta_5}$ as a consequence of the canonical two-stage
factorisation of $\Phi$.

\item \emph{CoHA incarnation.} $\bH_{\Delta_5}$ as a twisted
Cohomological Hall Algebra on the moduli of sheaves on $K3 \times E$,
with shuffle product in three spectral parameters truncated to the
CY slice $\sum \epsilon_i = 0$. Motivation: connects $\bH_{\Delta_5}$
to the Schiffmann--Vasserot framework and exhibits the
$Y^+_{\epsilon_1, \epsilon_2, \epsilon_3}$ positive-half presentation.

\item \emph{BPS incarnation.} $\bH_{\Delta_5}$ as the BPS algebra of
M-theory on $K3 \times E \times \RR^{4, 1}$, with root multiplicities
matching the $\frac{1}{4}$-BPS dyon counts weighted by $1/\Phi_{10}
= 1/\Delta_5^2$ (DVV 1997). Motivation: physical input makes the
mutliplicity generating function transparent and explains why
$\Delta_5^{-2}$ rather than $\Delta_5^{-1}$ appears as the natural
partition function.

\item \emph{Borcherds-double incarnation.} $\bH_{\Delta_5}$ as
$\mathrm{Heisdouble}(K3 \times E)$ --- the Drinfeld double of the
Mukai-lattice Heisenberg on $K3$ tensored with a fibre Heisenberg
on $E$, along the tangent-functor bimodule structure intrinsic to
a product CY. Motivation: exposes the $\mathsf{B}$-row witness
$K^{\kch} = 8$ via $\mathrm{ord}(S_{K3}^2 \otimes \tau_E) = 8$ as a
structural identity of Serre functors.

\item \emph{Quasi-categorical incarnation.} $\bH_{\Delta_5}$ as the
$E_1$-chiral algebra object in
$\mathrm{Alg}_{E_1}(D^b\mathrm{Coh}(K3 \times E)
^{\mathrm{MF-completed}})$, obtained by $\infty$-categorical
specialisation of the $E_3$-algebra $D^b\mathrm{Coh}(K3 \times E)$.
Motivation: the $(\infty, 1)$-categorical lane gives the cleanest
universal property; the chain-level presentation is then an
explicit $A_\infty$-model of this object.

\item \emph{Holographic incarnation.} $\bH_{\Delta_5}$ as the chiral
boundary algebra of the AdS$_3 \times S^3 \times K3$ near-horizon
geometry at the $N = 1$ CHL point, with $c_L^{\mathrm{reduced}} = 24$
(Brown--Henneaux, Maloney--Witten). Motivation: the gravity dual
exhibits $\bH_{\Delta_5}$ as a \emph{universal-family} boundary
algebra: every $\frac{1}{4}$-BPS state in the bulk is a state in
$\bH_{\Delta_5}$, and the partition function equals the Borcherds
denominator.
\end{enumerate}

\emph{Why seven candidates for equivalence.} The seven framings answer
seven different structural questions (``what is the denominator?'',
``which factorisation algebra?'', ``what Hall product?'', ``what BPS
count?'', ``what Serre-doubling?'', ``what $(\infty, 1)$-universal
object?'', ``what gravity dual?''), and the programme's Climax
Theorem is that the seven answers coincide. The pairwise equivalences
fall into three evidence-strength classes: (i)--(ii) follow from the
Stage-2 factorisation theorem (proved on the Koszul locus);
(iii)--(iv) are the CoHA-to-BPS bridge (Davison--Meinhardt 2015 at
the cohomological level); (v)--(vi) are an $(\infty, 1)$-categorical
Koszul-duality statement tracking the $\mathsf{B}$-row witness; (vii)
is the chiral-boundary holographic conjecture, contingent on a
rigorous proof of the $\mathrm{AdS}_3 / \mathrm{CFT}_2$ dictionary at
the level of chiral algebras on $K3 \times E$. All seven
incarnations agree on the rank-$3$ Cartan, the Borcherds weight $5$,
and the Euler-product coefficient data, giving seven independent
verification paths.
\end{theorem}

\textbf{Ghost}: without the seven-way inscription, the reader sees
disparate statements about $\bH_{\Delta_5}$ scattered across the
monograph (a factorisation algebra here, a BKM there, a BPS partition
function elsewhere) with no guidance on why these should agree.
\textbf{Error}: the original platonic ideals listed pairwise
equivalences but did not group them into an organised seven-way
compatibility statement. \textbf{Correction}: the seven framings are
inscribed with motivation lines making each incarnation's
mathematical role explicit; the pairwise evidence-strength structure
is declared so the reader can see which equivalences are theorems,
which are lifts, and which remain conjectural.

\subsection{Sibling specialisations}
\label{wn:subsec:healed-siblings-by-d}

\begin{conjecture}[$\Psi$-sibling structure, not $d$-stratified siblings]
\label{wn:conj:healed-siblings-psi}
The programme-canonical siblings under the $\Psi$-functor (Vol III \S11) are five:
\begin{itemize}
\item Monster: input $\mathrm{II}_{1, 1}$ sig $(1, 1)$, form $j(\sigma) - j(\tau)$ of weight $0$, rank $2$, $\ell_{\mathrm{Lusztig}} = 2$.
\item Conway super: input Leech sig $(24, 0)$, Conway denominator, $V^{s\natural}$-double, rank $2$ (super), on $\Psi^{\mathrm{metap}}$-branch (Scheithauer 2008; \emph{not} a fifth bosonic image, Wave 20 Kazhdan).
\item Enriques: input $E_8 \oplus \mathrm{II}_{1, 1}(2)$ sig $(1, 9)$, metaplectic $\Delta_5^{\mathrm{Enr}}$ weight $5/2$, $\mathbf H_{\Delta_5}/\!/\ZZ/2$, $\ell = 4$.
\item K3: input $\Lambda^{2, 1}_{II}$ sig $(2, 1)$, $\Delta_5$ weight $5$, rank $3$, $\ell = 8$.
\item Fake Monster: input $\mathrm{II}_{25, 1}$ sig $(25, 1)$, $\Phi_{12}$ weight $12$, rank $26$, $\ell = 50$.
\end{itemize}
Monstrous moonshine $V^\natural = \Sp_{T^{24}_{\mathrm{Leech}}/\ZZ_2, E_\tau}(\Phi^{\mathrm{FA}}_3(X_{\mathrm{Mon}}))$ via the Monster-specific Stage-1 on the Leech-quotient orbifold. The Hauptmodul property of McKay--Thompson series is Stage-1 uniqueness under $\mathrm{Co}_0$-equivariance, conditional on the Conway conjecture (Duncan--Mack-Ono 2015).
\end{conjecture}

\textbf{Ghost of retracted claim}: the original ``$d = 5$ cousin: Fake Monster ($K3 \times K3 \times E$, $\Sigma_4 = K3 \times K3$, $C = E$, $\kBKM(\Phi_{\mathrm{FM}}) = 12$)''. \textbf{Error}: Fake Monster has rank $26$ on $\mathrm{II}_{25, 1}$, not rank $h^{1, 1}(K3 \times K3 \times E) - 1 = 40$ or similar; the CY host $K3 \times K3 \times E$ does not produce the Fake-Monster lattice by any Stage-1 specialisation. \textbf{Correction}: the five $\Psi$-siblings are not $d$-stratified; they are $\Psi$-functor images on different lattice inputs (Vol III \S11). The dimension-stratified conjecture is retracted; the $\Psi$-sibling conjecture replaces it. Monster sits on its own Calabi--Yau-like orbifold $(T^{24}_{\mathrm{Leech}}/\ZZ_2) \times E$; K3-BKM sits on $K3 \times E$; Enriques sits on Enriques $\times E$; Fake-Monster has no compact CY host and lives on $\mathrm{II}_{25, 1}$ via $\Psi$-parallel transport.

\subsection{AdS$_3$ holography}
\label{wn:subsec:healed-AdS3}

\begin{theorem}[AdS$_3$ throat and CHL dyon counting]\ClaimStatusHealed
\label{wn:thm:healed-AdS3}
Near-horizon throat $\mathrm{AdS}_3 \times S^3/\ZZ_N \times (K3\text{-orbifold})$ with $c_N = 24 k_N$ and CHL relation $k_N = 24/(N + 1) - 2$ (Sen 2008). Explicit values: $(k_1, k_2, k_3, k_4, k_5, k_6) = (10, 6, 4, 3, 2, \text{no level-5 CHL; next is } 6 \text{ gives} 1)$. The dyon count
\[
d_N(Q, P) = \oint_{\mathcal C_N} \frac{e^{-i\pi (Q \cdot \Omega)}}{\Phi_{k_N}(\Omega)^2}\, d\rho\, d\sigma\, dv
\]
on the paramodular contour (Dijkgraaf--Verlinde--Verlinde 1997; Shih--Strominger--Yin 2005) has leading entropy $\log d_N = \pi \sqrt{N(Q^2)(P^2) - (Q \cdot P)^2} + (\ldots)$. Graviton finiteness equals $\mathrm{HH}^2_{E_2}$-vanishing (Wave 15 N1).
\end{theorem}

\textbf{Ghost}: the original ``$k_N = 24/(N+1) - 2$'' is correct, but the entropy formula ``$\log d_N = 2\pi \sqrt{\Delta/N} + (k_N + 2)\log\Delta$'' is garbled. \textbf{Error}: the Cardy-type leading entropy for CHL dyons is $\pi \sqrt{Q^2 P^2 - (Q \cdot P)^2 / N}$ with subleading $-(k_N + 2) \log (Q^2 P^2)$ (Sen 2008 eq. 4.2; Bhattacharya--Jatkar--Sen 2010). The $2\pi\sqrt{\Delta/N}$ form conflates with a different CFT quantity. \textbf{Correction}: the healed expression is the standard CHL entropy; the $\mathrm{HH}^2_{E_2}$-vanishing interpretation (Wave 15 N1) is retained.

\subsection{Adversarial retractions, healed}
\label{wn:subsec:healed-retractions}

\begin{theorem}[Retraction ledger, healed]
\label{wn:thm:healed-retractions}
\begin{enumerate}
\item $\widehat{\mathfrak{sl}}_3 \hookrightarrow \fg_{\Delta_5}$ from $\eta^9$ exponent: \ClaimStatusRetracted. True structure: real-root subalgebra is Feingold--Frenkel $F_3$ rank-3 hyperbolic Kac--Moody; $\eta^9$ in $\mathrm{Grit}(\eta^9 \vartheta_1) = \Delta_5$ is the Gritsenko lift's Jacobi-weight factor, not an exponent of any sub-Kac--Moody.

\item $\mathrm{VOA}[T_{24}] = L_{-6}(\mathfrak e_8)$: \ClaimStatusRetracted. Central-charge mismatch $c = -62 \neq -214$. True structure: $\mathcal V_{24} = H^0_{\mathrm{DS}}(L_{-2 + 1/22}(\mathfrak{sl}_2)^{\otimes 22})$ via iterated Drinfeld--Sokolov on the $\Sigma_{0, 24}$ pants decomposition (Arakawa 2017; Creutzig--Linshaw 2017). The central charge $-214 = (5 \cdot 24 - 13) \cdot (-1) / 6 \cdot 12 = 107/6 \cdot (-12)$ is $c_{2d} = -12 c_{4d}$ of $\mathcal T[A_1, \Sigma_{0, 24}]$ (Beem--Rastelli chiral-algebra map), not a product $22 \cdot 10 - 6$ (which does equal $214$ numerically but reads the wrong structural decomposition).

\item $\kBKM = \kch + \chi(\cO_{\mathrm{fiber}})$ universally: \ClaimStatusRetracted. The ``$N = 1$ coincidence'' that drove the conjecture is itself confabulated: at $N = 1$, $\kBKM(\Delta_5) = 5 \neq 0 = \kch(K3 \times E) + \chi(\cO_E) = 0 + 0$ (Vol III AP238, HEAL 2026-04-17). The universal law is Borcherds weight $\kBKM = c_N(0)/2$, which does \emph{not} decompose additively through $\kch$ and fibre structure. True structure: Theorem~\ref{wn:thm:healed-universal-kBKM}.

\item Fake Monster at $d = 3$ via Leech vs $h^{1, 1}(K3) = 20$: \ClaimStatusRetracted. See Remark~\ref{wn:rmk:healed-fake-monster}: the obstruction is rank $26$ vs compact-CY$_3$ $h^{1, 1} \leq 20$, but the Fake Monster is a $\Psi$-sibling on $\mathrm{II}_{25, 1}$, \emph{not} a compact-CY host shadow. The ``$d = 5$ resolution'' is also retracted; Fake Monster has no compact CY host.

\item $\chi_{\mathcal V_{24}}$ matches $\Delta_5^{-2}$ directly via Virasoro $(2, 45)$-minimal: \ClaimStatusRetracted. Three errors in the original Wave-15 N3 claim: arithmetic central-charge error ($-14432/121$ vs correct $-1312/11$); $\eta^{-48}$ coefficient divergence at $g_2$; Virasoro $(2, 45)$-minimal apparatus inapplicable (level matching fails). True structure: Theorem~\ref{wn:thm:healed-heis-mukai} at the Heisenberg level; the Virasoro-minimal match is a mirage.

\item Gaiotto curve $\Sigma_{2, 0}$: \ClaimStatusRetracted. True structure: $\Sigma_{0, 24}$ (24-punctured sphere) via Chacaltana--Distler 2010 Table 3 row 1: $c_{4d} = (5n - 13)/6 = 107/6$ at $n = 24$. The genus-$g$ closed-curve formula $c_{4d} = (12(g-1) + 7n)/6$ originally proposed fails already at $n = 4$ (gives $8/3$, not $7/6$; Wave 15 Gaiotto). The correct parent theory is $\mathcal T[A_1, \Sigma_{0, 24}]$ with $(n_v, n_h) = (63, 88)$, $c_{2d} = -214$.

\item $\Phi$ natively produces an $E_n$-chiral algebra on a curve: \ClaimStatusCorrected. True structure: two-stage (Theorem~\ref{wn:thm:healed-two-stage}); Stage~1 produces $E_d$-hFA on $X$, Stage~2 specialises to a curve.

\item Shifted-symplectic table terminates at $d = 4$: \ClaimStatusCorrected. True structure: PTVV admits arbitrary integer shifts; shift law $\mathrm{shift} = d - 4$ (CPTVV 2017). The $E_n$-label on bulk quantum observables on flat $\CC^d$ is $E_d$, not ``$E_{d - 0}$''.

\item Hecke trace $\lambda_p(\Delta_5)$ equals $1$: \ClaimStatusCorrected. True structure: $\Delta_5$ lives on spin cover; $\Delta_5^2 = \Delta_{10}$ with $\lambda_p(\Delta_{10}) = a_p(\Delta_{E_6}) + p^8 + p^9$ (Andrianov--Ikeda; Saito--Kurokawa CAP; Vol I TRUTH I.7). The ``multiplicity-one = Hecke eigenvalue'' conflation is the Wave 14 K1 error.

\item Twisted-twined $H^3(C_{\Mtf}(g_N); U(1)) = 0$ uniformly: \ClaimStatusCorrected. True structure: non-uniform across admissible cells; $2A$ carries $\ZZ/2$, $2B$ carries $\ZZ/4$ residual discrete torsion (Milgram 1995; Benson 1998). The $\kBKM(\Phi_2) = c_2(0)/2$ is cocycle-free at the weight-index level; the $M_{24}$-torsion dresses the paramodular multiplier, not the Borcherds weight.
\end{enumerate}
\end{theorem}

\subsection{Residual frontier}
\label{wn:subsec:healed-frontier}

\begin{itemize}
\item $\kBKM$ cross-volume AP5 lock: the $(5, 12)$ dual-reading is preserved in the canonical preamble row 57 (Vol III) / canonical row 57 compatible-dual-reading 3; each inscription must name its input denominator (Pattern 411).
\item Elliptic-surface $(\Sigma_2, C) = (\mathcal E, \PP^1)$: live structural frontier (Wave~16 S2).
\item Non-CHL $N = 7$: order-$4$ gerbe Cheeger--Simons realisation open.
\item Fake-Monster explicit $\Psi$-image generators beyond rank-26 RTT skeleton.
\item Non-simply-laced $B/C/F/G$ extension of the all-orders Yangian compound (via folding; Wave 16 T1).
\item Lattice-polarised rank-$\geq 3$ $\fg_L$ family (Wave~16 U2 extension).
\item $\Phi_4$ framework gap: Kapustin--Rozansky--Saulina 3d/4d dichotomy blocks the CY-$4$ extension.
\end{itemize}

\subsection{The healed picture in one line}
\label{wn:subsec:healed-one-line}

\emph{Five $\Psi$-siblings on five input lattices (Monster / Conway-super / Enriques / K3 / Fake-Monster) realise the BKM crown from distinct automorphic-product data, not five-fold $(\Sigma_{d-1}, C)$-shadows of a single CY$_d$; the two-stage factorisation $\Phi_d = \Sp_{\Sigma_{d-1}, C} \circ \Phi^{\mathrm{FA}}_d$ applies per Calabi--Yau host; the universal $\kBKM = c_N(0)/2$ Borcherds-weight rule is a single formula whose concrete value indexes the input denominator ($\Delta_5 \mapsto 5$, $\Phi_{10} \mapsto 5$, $\Phi_{12} \mapsto 12$); on $K3 \times E$ the canonical Stage-1 $\cF_{K3 \times E}$ has rank $24$ Mukai-grading orthogonal to the rank $3$ BKM Cartan; the directionality ``canonical factorisation algebra per CY, specialisation by transverse cycle and curve'' survives, while the ``many-from-one'' unification is realised by the $\Psi$-functor (operadic sibling structure, Wave 20 Gaiotto) and not by Stage-2 shadows of a shared $E_3$-hFA across hosts.}
